\begin{table}[t]
\centering
\caption{Standardized Effect Sizes}
\label{tab:sde}
\begin{adjustbox}{max width=\textwidth}
\begin{tabular}{lcccccc}
\hline\hline
Outcome & $\hat{\beta}$ & SE & SD($Y$) & SDE & SE(SDE) & Classification \\
\hline
\multicolumn{7}{l}{\textit{Panel A: Pooled}} \\[3pt]
Total casualties & 0.0777 & 0.0236 & 9.633 & 0.0081 & 0.0024 & Small positive \\
Total injuries & 0.0757 & 0.0222 & 8.554 & 0.0089 & 0.0026 & Small positive \\
Total deaths & 0.0020 & 0.0020 & 1.132 & 0.0018 & 0.0018 & Null \\
Any casualty & 0.0016 & 0.0012 & 0.238 & 0.0069 & 0.0048 & Small positive \\
\\
\multicolumn{7}{l}{\textit{Panel B: Heterogeneous}} \\[3pt]
Casualties, EF2+ tornadoes & 1.0159 & 0.5857 & 28.702 & 0.0354 & 0.0204 & Small positive \\
Casualties, high-mobile-home counties & 0.1450 & 0.0559 & 13.280 & 0.0109 & 0.0042 & Small positive \\
\hline\hline
\end{tabular}
\end{adjustbox}
\begin{tablenotes}[flushleft]\footnotesize
\item \textit{Notes:} \textbf{Country:} United States. \textbf{Research question:} Does the quality of tornado warnings issued by NWS Weather Forecast Offices, as measured by average lead time, causally affect tornado casualties in adjacent counties assigned to different offices? \textbf{Policy mechanism:} The 122 NWS Weather Forecast Offices were assigned fixed County Warning Areas during the 1990s modernization following administrative convenience rather than tornado risk, creating persistent differences in warning performance across arbitrary boundaries. \textbf{Outcome definition:} Total casualties (injuries plus deaths) per tornado event from the SPC Storm Data. \textbf{Treatment:} Continuous; WFO-level average lead time in minutes between warning issuance and tornado touchdown, averaged over 2008--2024 from IEM verification data. \textbf{Data:} SPC tornado records (2008--2024), IEM Cow verification data, NWS CWA shapefiles, and Census ACS; 21,346 tornado-event-by-boundary-pair observations across 1,602 pairs and 106 WFOs. \textbf{Method:} OLS with boundary-pair and year fixed effects, two-way clustering by WFO and year. \textbf{Sample:} US tornado events in counties adjacent to a WFO boundary, restricted to county-level records with valid FIPS and coordinates. SDE $= \hat{\beta} / \text{SD}(Y)$ where SD($Y$) is the full-sample standard deviation of the outcome. Classification refers to magnitude, not statistical significance: Large ($|$SDE$| > 0.15$), Moderate ($0.05$--$0.15$), Small ($0.005$--$0.05$), Null ($< 0.005$).
\end{tablenotes}
\end{table}
